\clearpage
\section{Verification record}
\label{sec:I:verification}
Appendix~\ref{app:I:result-index} identifies every source result and its
printed statement. A recorded independent proof reading counts as a check;
a producing run's status does not. We compared the retained source arguments
with the printed proofs, including the moved re-derivations and all distinct
methods. An exact-arithmetic rerun checks its stated algebraic assertions;
a ball-arithmetic rerun checks its inequalities and explicit analytic tails.
Neither is treated as a check of unrelated geometric steps. Source checks
match the theorem actually used, its edition, and all its hypotheses.

The current full readings, derivations, original-source checks and
certificate reruns were performed by GPT-6 Astra (OpenAI, ChatGPT Pro mode,
effort Pro), 17 November 2026. The per-result record preserves the earlier
readings by Claude Fable 5.2 (Anthropic), effort high or effort max as recorded,
12--16 November 2026, and Claude Opus 5.5 (Anthropic), effort max,
12--16 November 2026. The complete machine-readable verification record is
included with the programs and logs; proofs are printed in this paper.

The citations to Bertrand's algebraicity theorem, the original
Tate--Silverman specialization sources, Rubin's CM work and the sigma-value
survey are retained only at the expressly qualified comparisons. Their
unverified entries support no implication in a proof. The rejected
algebraicity implication, the absence of a general unconditional norm
transcendence argument and the distinction between an exact leading factor
and a full Birch--Swinnerton-Dyer assertion are stated at their points of use.

\section{Provenance}
The working edition was AI-produced by GPT-5.6 Sol (OpenAI) in Codex at
reasoning effort ultra, 24 August--9 September 2026, with a few inputs from
GPT-5.6 Sol (OpenAI, ChatGPT Pro mode), 24 August--9 September 2026.
The course continuation records GPT-6 Astra (OpenAI), Codex reasoning
effort ultra, October 2026; its chapter dates and retained prerequisite
provenance are distinguished in the accompanying record.
The retained chain-level prerequisite was AI-produced by GPT-6.1 Sol
(OpenAI), reasoning effort ultra, October 2026, and selected and supplemented
by GPT-6 Astra (OpenAI), reasoning effort ultra, October 2026.
The unchanged homotopy-fibres provider names GPT-6.1 Sol (OpenAI),
reasoning effort ultra, with date not recorded; the receiving path-fibre
companion separately names GPT-6 Astra (OpenAI), reasoning effort ultra,
October 2026. These provenance dates are not conflated.
The six research runs were AI-produced by GPT-6 Astra (OpenAI, ChatGPT Pro
mode), 12--16 November 2026. This writing, the additional readings and the
comparison with the printed proofs were produced by GPT-6 Astra (OpenAI,
ChatGPT Pro mode, effort Pro), 17 November 2026.
The recorded readings were by Claude Fable 5.2 (Anthropic), effort high and
effort max, 12--16 November 2026, and Claude Opus 5.5 (Anthropic), effort max,
12--16 November 2026. Exact source-to-print label changes, duplicated
statements, moved proofs, repaired gaps and rejected assertions are recorded
in the concordance and accompanying verification record.

\paragraph{Effort terminology.}
The recorded effort settings describe computational modes and do not certify
the mathematics. Anthropic's official effort documentation identifies
\texttt{max} as its maximum capability setting
\cite[How effort works; Effort levels]{AnthropicEffort2026}.
We retain the recorded ChatGPT Pro-mode designation.
OpenAI's official documentation distinguishes reasoning mode from reasoning
effort and describes Pro mode as using additional computation
\cite[Reasoning mode]{OpenAIReasoning2026}.
It does not identify Pro mode with one universal maximum effort value; we
therefore do not make that identification.

\section{Declarations}
This paper is AI-produced. No human has checked its mathematics, and no
human understanding of any part is claimed. We name no person as an author
of or contributor to this work. The authors of other works receive credit
for their own results. No formal proof-assistant certification is claimed.
Our text is dedicated to the public domain under CC0 1.0.

\paragraph{Data statement.}
The programs, certificates, data and run records are held at
\url{https://github.com/KokunoYumeto/s6-complex-structure-workbench}.
The companion archive contains the programs written or rerun for this
edition, their commands and outputs, the result inventory, the notation
dictionary and the correction record. Appendix~\ref{app:I:programs}
specifies the programs' roles and how to run them.

\bibliographystyle{apalike}
\bibliography{S6_09_PAPERS_references}
